\section*{Anexo 4.1-M --- Verificación de continuidad lógica}
\phantomsection
\label{anx:M}
\addcontentsline{toc}{section}{Anexo 4.1-M --- Verificación de continuidad lógica}
Este anexo especifica pruebas de aceptación, no resultados de ensayos. Cada ejecución debe conservar versión de software y configuración, datos de carga, relojes sincronizados, comandos de inyección de falla, logs con correlación, resultados esperados y observados, incidencias y acta. Sin estos artefactos no se acredita cumplimiento. Los responsables indicados son roles propuestos de ejecución y aprobación, no firmas ya obtenidas.

\subsection*{AL-DTE-01. Salida de 96 camiones y falla del emisor}
La prueba verifica el control de M5/ACL/ERP frente a la ventana de 05:30--07:00 (Caso 02, cap.~10 y cap.~14). Operaciones aporta el programa de 96 salidas, la relación de cargas y la cantidad real de guías por camión. No se supone una guía por camión: la carga de emisión se deriva de los destinos y documentos necesarios. Tributación del CLIENTE y el responsable del ERP deben identificar por escrito el procedimiento permitido, sus límites, autorizaciones, recuperación y documentos de respaldo.

Se ensayan cuatro escenarios: indisponibilidad del ERP antes de preparar la carga; guía emitida cuya respuesta se pierde; falla del canal tributario después de cerrar carga; y cambio de carga cuando ya existe guía. En cada uno se verifica que la misma clave no genere dos documentos, que no se libere una carga distinta de la documentada y que los camiones no afectados continúen. El procedimiento autorizado, si aplica, registra responsable, causal, hora, documento y conciliación posterior. La aplicación no crea un emisor sustituto ni habilita un botón genérico para saltar el control.

El criterio de aceptación conjunta es disponer de documentación válida para las 96 salidas programadas, sin interrupción atribuible a la solución dentro de la ventana y sin duplicados ni emisión retroactiva desde terreno. Si la única respuesta posible es retener camiones, el ensayo preserva el control tributario, pero falla el criterio de continuidad. La anticipación de guías debe probarse con cambios tardíos y no se acepta como única mitigación. El responsable de Arquitectura debe elevar el conflicto al CLIENTE con el resultado del ensayo y la alternativa de continuidad del ERP; el requisito de indisponibilidad cero permanece vigente.

\subsection*{AL-DR-01. Pérdida de sitio durante una caída WAN}
RT-07.04 mantiene RTO máximo de 4 horas y RPO máximo de 15 minutos para servicios críticos. El punto de recuperación se mide con la última transacción de negocio confirmada y recuperable, no con el último mensaje enviado ni con la existencia de un respaldo. Para cada flujo se registra la marca de agua de confirmación durable fuera del dominio de falla local. Se define $RPO_{observado}=t_{incidente}-t_{ultimo\_punto\_consistente}$, verificando además las operaciones efectivamente perdidas por UUID y secuencia.

El ensayo confirma primero el caso con enlace disponible. Después interrumpe todos los caminos WAN del sitio, mantiene carga local durante al menos 60 minutos y simula la pérdida completa del origen sin reutilizar sus discos, colas ni WAL. El equipo recupera desde el sitio secundario y compara pedidos, stock, preparación, bloqueos sanitarios y evidencias contra un oráculo de prueba externo. Se repite con una desconexión de 24 horas para comprobar el límite de la envolvente offline. Se cronometran RTO y RPO reales, no tiempos del proceso de copia por separado.

La protección requerida es una copia durable de transacciones y evidencias críticas fuera del dominio de falla, con un desfase no mayor de 15 minutos aun en el escenario acordado. Un segundo medio de comunicación solo sirve si permanece disponible ante la falla ensayada; dos enlaces que caen juntos no la proporcionan. DMS, WAL y RabbitMQ almacenados únicamente en el sitio perdido tampoco sirven. Con aislamiento total prolongado y sin extracción externa, el diseño descrito no puede demostrar ese RPO. Suspender escrituras al minuto 15 protegería la antigüedad de los datos, pero contradice la autonomía local de 24 horas y no se adopta como solución implícita.

La brecha AL-BR-01 queda identificada para decisión de Arquitectura, Infraestructura y CLIENTE: diseñar y presupuestar una protección externa independiente y probarla, o tramitar una resolución formal del mandante sobre el escenario. Su registro no equivale a elevar una consulta ya enviada ni a obtener dispensa. La aceptación no se cierra mientras persista un resultado superior a 15 minutos o no exista una resolución formal aplicable; este apartado no modifica el diseño físico ni rebaja el requisito.

\subsection*{AL-OFF-01. Autonomía de 24 horas con dos relevos}
La prueba cubre RT-03.10 y RT-03.13, con los perfiles reales de Talca, Concepción y un cross-dock; los otros sitios deben ejecutar el mismo protocolo antes de su ola. Se preinscriben usuarios, dispositivos, turnos y permisos. El corte comienza justo antes de la renovación horaria del manifiesto de 26 horas, con caché de identidad de 24 horas y reloj local controlado. Se bloquea tanto la WAN como el IdP durante 24 horas; no basta cerrar una sesión del navegador.

A las 8 y 16 horas, personas distintas inician turno con PIN personal en terminal enrolado. Se verifican firma, vigencia, rol, sitio y turno; se rechazan usuario no preinscrito, manifiesto alterado, PIN repetidamente erróneo, dispositivo no autorizado y privilegio administrativo nuevo. Se comprueba que expirar la caché descriptiva no renueve permisos ni invalide indebidamente una asignación firmada vigente; el verificador debe disponer en el manifiesto de los atributos mínimos para decidir. Se reinicia un componente local y se verifica persistencia de sesiones autorizadas, outbox y bitácora sin reutilizar claves personales.

Durante el corte se ejecutan recepción, reserva local, picking, bloqueos térmicos, despacho con documento disponible y captura de incidencias. Se incluye el intento de nueva salida sin guía: su aceptación depende de AL-DTE-01, no de aprobar el PIN. Tras reconectar se drenan colas con reenvíos, duplicados y eventos fuera de orden; se comparan stock, pedidos, POD, permisos y auditoría con el oráculo. Una revocación central desconocida durante el corte se aplica al reconectar y sus acciones del intervalo se revisan; no se declara revocación instantánea offline.

Se acepta cuando ambos relevos funcionan con personas identificadas, no hay autorización indebida, todas las operaciones confirmadas reaparecen una sola vez en el destino y cada conflicto conserva regla y responsable. El informe registra latencias, ocupación máxima, antigüedad de cola y tiempo de drenaje bajo carga de septiembre. Seguridad, Calidad y Operaciones firman los resultados. La autonomía de despacho no se declara completa si queda abierto AL-DTE-01.
